\abschnitt{Vision: FabCities}
% Maßstäbe von zu Hause bis zum Planeten. Rechts zwei Balken: Material
% bleibt in der Nähe (bis zur Stadt), Daten gehen um die Welt.
\vspace{6mm}
\begin{tikzpicture}[x=1mm,y=1mm]
  \foreach \i/\r/\name/\weite in {0/2.5/Home/,1/4.5/Neighbourhood/1 km,
      2/6.5/City/10 km,3/8.5/Region/{1,000 km},4/10.5/Planet/{10,000 km}}{
    \draw[faublau,line width=1.25mm] (10.5,-\i*34) circle (\r);
    \node[anchor=base west,inner sep=0,font=\beschriftungschrift\bfseries,
      text=faublau] at (28,-\i*34-4) {\name};
    \node[anchor=base west,inner sep=0,font=\beschriftungschrift,
      text=fautech] at (28,-\i*34-17) {\weite};
  }
  % Material: lokal
  \draw[faublau,line width=2.5mm,line cap=round] (185,4) -- (185,-72);
  % Daten: global
  \draw[fautech,line width=2.5mm,line cap=round] (208,4) -- (208,-140);
  % Legende
  \draw[faublau,line width=2.5mm,line cap=round] (4,-172) -- (17,-172);
  \node[anchor=base west,inner sep=0,font=\beschriftungschrift]
    at (28,-176) {materials stay local};
  \draw[fautech,line width=2.5mm,line cap=round] (4,-190) -- (17,-190);
  \node[anchor=base west,inner sep=0,font=\beschriftungschrift]
    at (28,-194) {data travels globally};
\end{tikzpicture}
\par\vspace{12mm}
Global exchange of data, with local harvesting and recycling of
materials.
